\section{Análisis}

A 1\,kHz la separación entre tramas es suficiente y casi todas las respuestas
traen el encabezado correcto, pero el 84\,\% falla el CRC. En el mapa de
bits, las posiciones con más errores son el bit 7 de los bytes 1, 2 y 7
(alrededor de 40\,000 cada una, contra 16\,000 como máximo en cualquier otra). Al compararlos con la respuesta esperada, en el 77\,\% de los
151\,782 bits 7 erróneos el valor recibido es el del bit 0 del byte anterior,
que es lo último que había en SDO. Por ejemplo, la respuesta
esperada \texttt{5A 01 00 00 00 00 00 EC} llega como
\texttt{5A 01 00 00 00 00 00 6C}: el bit 7 del CRC toma el valor del bit 0
del status (0).

Esto indica que el dsPIC no alcanza a poner en SDO el primer bit de cada
byte antes de que el FT2232H lo muestree. En modo 0 el maestro lee ese bit
en el primer flanco de subida, medio periodo de SCK (500\,ns a 1\,MHz)
después de terminar el byte anterior, y el módulo SPI del dsPIC, sincronizado
con un reloj de periférico de 4\,MHz (250\,ns por ciclo), necesita del orden
de ese tiempo para cargar el byte siguiente en el registro de corrimiento.
Las pruebas con SCK de 500 y 250\,kHz lo confirman: los errores de CRC bajan
del 84\,\% al 0.5\,\%.

Además, alrededor del 6\,\% de las respuestas trae número de secuencia 0 y
posición 0, que es la respuesta que el firmware construye al arrancar. En
compilación de producción, los manejadores de trampas generados por MCC
(\texttt{traps.c}) ejecutan \texttt{reset}; entre ellas está el error de
dirección del DMA. Estas respuestas aparecen también con SCK más lento y en
rachas de tramas consecutivas, por lo que no se descarta que se trate del DMA
retransmitiendo la respuesta inicial en lugar de un reinicio. Debe
verificarse con una compilación de depuración (que se detiene en la trampa) o
leyendo \texttt{RCON} al arrancar.

\subsection*{Recomendaciones}

\begin{enumerate}
\item Elevar $F_{cy}$ con el PLL. Acorta la interrupción de fin de trama
  (necesario para la prueba a máxima velocidad) y la latencia del SPI
  esclavo, que causa los errores de CRC.
\item Mientras tanto, bajar SCK a 500\,kHz elimina casi todos los errores de
  CRC; la transferencia pasa de unos 130 a unos 190\,$\mu$s, todavía muy por
  debajo del periodo de 1\,ms.
\item Determinar la causa de las respuestas de arranque antes de usar el
  bloque en lazo cerrado.
\end{enumerate}
